% Copy this file into a release's tex/ directory and edit the prose only after
% checking the claim-to-artifact matrix and the candidate scientific status.
\subsection{Directional activation energies}
\label{sec:directional-activation}

For a fixed NEB energy trace with images ordered from endpoint 1 to endpoint
10, we report both directions. The forward value is
\(E^{\rightarrow}_{\mathrm{act}}=\max_i(E_i-E_1)\), while the reverse value
is \(E^{\leftarrow}_{\mathrm{act}}=\max_i(E_i-E_{10})\). Thus, a downhill
forward trace can still have a nonzero reverse activation energy. These two
numbers are not two different calculations; they are directional summaries of
the same fixed path.

Figure~\ref{fig:gap2ql-bidirectional} and Table~\ref{tab:gap2ql-bidirectional}
show the paired forward and reverse values for the gap-to-QL comparison. The associated table and geometry audit
must be included in the evidence package before this paragraph is promoted from
`DRAFT` to `SUPPORTED`.

\begin{figure}[t]
  \centering
  \includegraphics[width=\linewidth]{figures/figure_gap2ql_bidirectional_activation.pdf}
  \caption{Forward and reverse activation energies extracted from the same
  fixed NEB traces. Replace the caption with the release-specific scientific
  status and cite the corresponding table and claim matrix row.}
  \label{fig:gap2ql-bidirectional}
\end{figure}

\begin{table}[t]
  \centering
  \begin{tabular}{ll}
    Direction & Definition \\
    Forward & $\max_i(E_i-E_1)$ \\
    Reverse & $\max_i(E_i-E_{10})$ \\
  \end{tabular}
  \caption{Definitions used by the bidirectional activation-energy summary. Replace with the release-specific table and cite its source CSV.}
  \label{tab:gap2ql-bidirectional}
\end{table}
